\documentclass[11pt]{amsart}
\usepackage{amsmath,amssymb,amsthm,booktabs}
\usepackage[margin=1.15in]{geometry}
\usepackage[hidelinks]{hyperref}
\newtheorem{theorem}{Theorem}[section]
\newtheorem{proposition}[theorem]{Proposition}
\newtheorem{lemma}[theorem]{Lemma}
\theoremstyle{remark}
\newtheorem{remark}[theorem]{Remark}
\newcommand{\PP}{\mathbb{P}}
\newcommand{\EE}{\mathbb{E}}
\newcommand{\one}{\mathbf{1}}
\newcommand{\doi}[1]{\href{https://doi.org/#1}{doi:\nolinkurl{#1}}}
\newcommand{\arxiv}[1]{\href{https://arxiv.org/abs/#1}{arXiv:\nolinkurl{#1}}}
\setlength{\emergencystretch}{2.5em}

\title[Subcritical components on near-regular graphs]{Polynomial components in spectrally subcritical percolation on
near-regular graphs}
\author{Yitzchak Shmalo}
\address{Einstein Institute of Mathematics, The Hebrew University of Jerusalem, Jerusalem, Israel}
\email{yspennstate@gmail.com}
\date{November 2026}
\subjclass[2020]{05C80, 60K35, 05C50}
\keywords{Bernoulli percolation, finite graphs, spectral radius, largest component, branching processes}

\raggedbottom
\begin{document}
\begin{abstract}
Let $G$ be a finite graph with spectral radius $\lambda$ and maximum degree $\Delta$, and keep each edge independently
with probability $p=(1-\varepsilon)/\lambda$. Christoph, Morawski and Wigderson asked for the typical size of the largest
component in this regime, singling out graphs that are close to regular in the sense $\Delta\le(1+\delta)\lambda$. We
show that the answer depends on the product $(1-\varepsilon)(1+\delta)$. When it is less than one, a known
maximum-degree argument gives components of logarithmic order. When it is greater than one, there are connected graphs
with $\Delta\le(1+\delta)\lambda$ whose largest component has at least $n^c$ vertices with high probability, for a
constant $c>0$, and the spectral radius can be taken of order $n^a$ for any $0\le a<1$. When it equals one, the largest
component has $O((\log n)^3)$ vertices with high probability. We also sharpen the resolvent bound of Christoph, Morawski
and Wigderson for graphs with bounded ratio $\Delta/\lambda$. The optimal order in the polynomial regime remains open.
\end{abstract}
\maketitle

\section{Introduction}\label{sec:intro}

Let $G$ be a finite simple graph on $n$ vertices with adjacency matrix $A$, spectral radius $\lambda=\lambda(G)$ and
maximum degree $\Delta=\Delta(G)$. Write $G_p$ for the random subgraph that keeps each edge independently with
probability $p$, and $L_1(G_p)$ for the number of vertices of its largest component. Percolation on general finite
graphs has a long history; see Alon, Benjamini and Stacey~\cite{ABS04} for the role of expansion, and
Bollobás, Borgs, Chayes and Riordan~\cite{BBCR10} for dense graph sequences, where the threshold is $1/\lambda$.

Christoph, Morawski and Wigderson~\cite{CMW26} proved that $1/\lambda$ is the critical probability for every graph in a
strong sense. In the subcritical regime $p\le(1-\varepsilon)/\lambda$ they showed~\cite[Theorem~1.4]{CMW26} that
\begin{equation}\label{eq:cmw}
L_1(G_p)=O_\varepsilon\Bigl(\frac{\sqrt{n\Delta}}{\lambda}\log n\Bigr)\qquad\text{with high probability},
\end{equation}
and they asked~\cite[Question~5.2]{CMW26} for the tight upper bound in terms of $\Delta$ and $\lambda$. They single out
graphs that are close to regular, $\Delta\le(1+\delta)\lambda$, with $\delta$ allowed to depend on $n$, and they note
that finite $d$-ary trees have polynomial components~\cite[p.~20, footnote~7]{CMW26}. For trees the ratio
$\Delta/\lambda$ is about $\sqrt d/2$, far from one.

Throughout, $\varepsilon\in(0,1)$ is fixed and $q=1-\varepsilon$. The answer for near-regular graphs depends on whether
$q(1+\delta)$ is below, above or equal to one.

\begin{theorem}\label{thm:lower}
Let $\delta>0$ satisfy $q(1+\delta)>1$ and let $0\le a<1$. There are connected graphs $G_h$, $h\ge1$, with
$n_h\to\infty$ vertices, $\Delta(G_h)\le(1+\delta)\lambda(G_h)$ and $\lambda(G_h)=\Theta(n_h^{a})$, and a constant
$c>0$ depending on $\varepsilon$, $\delta$ and $a$, such that for $p_h=q/\lambda(G_h)$,
\[
\PP\bigl(L_1((G_h)_{p_h})\ge n_h^{c}\bigr)\to1.
\]
\end{theorem}

\begin{proposition}\label{prop:strict}
If $\Delta\le(1+\delta)\lambda$ with $q(1+\delta)<1$ and $p\le q/\lambda$, then $L_1(G_p)=O_{\varepsilon,\delta}(\log
n)$ with high probability. The order $\log n$ is attained by the cycle $C_n$.
\end{proposition}

\begin{theorem}\label{thm:boundary}
If $p\lambda\le q$ and $p\Delta\le1$, then $L_1(G_p)=O_\varepsilon((\log n)^3)$ with high probability. In particular
this holds when $\Delta\le(1+\delta)\lambda$ with $q(1+\delta)=1$ and $p\le q/\lambda$.
\end{theorem}

Proposition~\ref{prop:strict} is a special case of a theorem of Diskin, Erde, Kang and
Krivelevich~\cite[Theorem~2]{DEKK25}, which needs only $p\Delta\le1-\kappa$; we include the short proof to fix
constants. Theorem~\ref{thm:boundary} uses a multitype branching process dominated through the resolvent vector, in the
spirit of the subcritical analysis of~\cite[Theorem~3]{BBCR10}.

The constructions behind Theorem~\ref{thm:lower} join a clique, which sets the spectral radius, to a regular tree of
slightly larger degree, which supplies a supercritical Galton--Watson process at $p=q/\lambda$. The tree alone has small
spectral radius, so it does not raise $\lambda$, and the clique has degree close to $\lambda$, so the ratio
$\Delta/\lambda$ stays below $1+\delta$. For a fixed ratio $\Delta/\lambda$ close to one, Theorem~\ref{thm:lower} needs
$\varepsilon<\delta/(1+\delta)$; for $\delta=o(1)$ and fixed $\varepsilon$ Proposition~\ref{prop:strict} applies
eventually. So the three statements describe the near-regular case for every fixed pair $(\varepsilon,\delta)$ up to
the exponent in the polynomial regime and the power of $\log n$ at the boundary.

The bound~\eqref{eq:cmw} loses in two places: a walk count from a vertex is bounded by $\sqrt{d(x)n}\,\lambda^{k-1}$, and
the resulting weights are compared globally. Keeping the trivial bound $\Delta^k$ beside it gives the following
refinement. Put
\[
R=\frac{\Delta}{\lambda},\qquad B=\frac{\sqrt{n\Delta}}{\lambda}.
\]

\begin{theorem}\label{thm:refine}
Fix $a$ with $1<a<1/q$, let $\theta=1-qa$ and
\[
U_a=1+\sum_{j\ge1}a^{-j}\min\{R^j,B\}.
\]
If $n\ge2$ and $p\le q/\lambda$, then
\[
\PP\Bigl(L_1(G_p)\ge\big\lceil4\theta^{-2}U_a(1+3\log n)\big\rceil\Bigr)\le n^{-2}.
\]
\end{theorem}

For fixed $R>1/q$ this gives (Remark~\ref{rem:exponent}), for every $0<\zeta<-\log q/\log R$,
\begin{equation}\label{eq:refine}
L_1(G_p)=O_{\varepsilon,R,\zeta}\bigl(B^{\log(qR)/\log R+\zeta}\log n\bigr)\qquad\text{with high probability},
\end{equation}
and the exponent of $B$ is below one, which improves~\eqref{eq:cmw} along sequences with $B\to\infty$. The upper
exponent in~\eqref{eq:refine} and the lower exponent in Theorem~\ref{thm:lower} differ, and we do not know the optimal
one. Question~5.2 of~\cite{CMW26} asks for a bound for every triple $(n,\Delta,\lambda)$, and the results here do not
give it.

Section~\ref{sec:lower} proves Theorem~\ref{thm:lower}, Section~\ref{sec:strict} Proposition~\ref{prop:strict},
Section~\ref{sec:boundary} Theorem~\ref{thm:boundary} and Section~\ref{sec:refine} Theorem~\ref{thm:refine} with the
exponent~\eqref{eq:refine}. Section~\ref{sec:examples} compares the bounds on small examples and describes the
simulations.

\paragraph{Conventions}
Logarithms are natural. Asymptotic statements refer to $n\to\infty$, and ``with high probability'' means with
probability tending to one. The constants implied by $O(\cdot)$ and $\Theta(\cdot)$ depend only on the parameters
written as subscripts, or in Theorem~\ref{thm:lower} on the fixed construction. We write $\one$ for the all-ones
vector and $e_x$ for the standard basis vector at $x$. The feasible parameters satisfy $\sqrt\Delta\le\lambda\le\Delta\le n-1$: the first inequality
follows from $(A^2)_{xx}=d(x)$ and the second from the row sums of $A$. We write $W_k(x)=(A^k\one)_x$ for the number of
walks of length $k$ from $x$. For $k\ge1$, by Cauchy--Schwarz,
\begin{equation}\label{eq:walk}
W_k(x)=\langle Ae_x,A^{k-1}\one\rangle\le\sqrt{d(x)}\,\lambda^{k-1}\sqrt n,
\end{equation}
and trivially $W_k(x)\le\Delta^k$. A symbol table is in Appendix~\ref{app:notation}.

\section{Polynomial components above the margin}\label{sec:lower}

This section builds the graphs of Theorem~\ref{thm:lower} and finds a supercritical Galton--Watson tree inside their
percolation clusters.

\subsection{The construction}
Let $k,b\ge2$ be integers. For $h\ge1$ let $T_{b,h}$ be the complete rooted $b$-ary tree with levels $0,\dots,h$, and
let $H_h$ be the graph obtained from the disjoint union of $T_{b,h}$ and the clique $K_{k+1}$ by adding one edge, the
bridge, between the root of the tree and one vertex of the clique. For $m\ge1$, let $H_h[m]$ be the graph obtained by
replacing each vertex of $H_h$ by an independent set of $m$ vertices, its block, and each edge by the complete bipartite
graph between the two blocks. Its adjacency matrix is the Kronecker product $A_{H_h}\otimes J_m$, where $J_m$ is the
$m\times m$ all-ones matrix.

\begin{lemma}\label{lem:spectral}
If $k\le b$ and $k\ge2\sqrt b$, then $k\le\lambda(H_h)\le k+1$ and $\Delta(H_h)=b+1$. Consequently
$\lambda(H_h[m])=m\lambda(H_h)$ and $\Delta(H_h[m])=m(b+1)$.
\end{lemma}

\begin{proof}
For a real vector $x$ on the tree, using $2|x_ux_v|\le b^{-1/2}x_u^2+b^{1/2}x_v^2$ for each parent $u$ of $v$,
\[
|x^{\mathsf T}A_{T_{b,h}}x|\le\sum_{u\text{ parent of }v}\bigl(b^{-1/2}x_u^2+b^{1/2}x_v^2\bigr)\le2\sqrt b\,\|x\|_2^2,
\]
since each vertex has at most $b$ children and at most one parent. So $\lambda(T_{b,h})\le2\sqrt b\le k$, and the
disjoint union of the tree and $K_{k+1}$ has spectral radius $k$. The bridge has adjacency matrix of norm one, so
$\lambda(H_h)\le k+1$, and $\lambda(H_h)\ge\lambda(K_{k+1})=k$ by monotonicity of the spectral radius. In $H_h$ an
internal tree vertex has degree $b+1$, the root has $b$ children and the bridge, the clique vertex on the bridge has
degree $k+1\le b+1$, and all other degrees are smaller. The eigenvalues of $A_{H_h}\otimes J_m$ are $m$ times those of
$A_{H_h}$, and degrees are multiplied by $m$.
\end{proof}

\subsection{Choice of parameters}
Suppose $q(1+\delta)>1$ and fix $r$ with $1/q<r<1+\delta$. For an integer $k$ put $b=\lfloor rk\rfloor$ and
\[
\mu=b\bigl(1-e^{-q/(k+1)}\bigr).
\]
As $k\to\infty$ we have $b/k\to r$ and $\mu\to rq>1$. Fix $k$ so large that
\begin{equation}\label{eq:params}
2\sqrt b\le k\le b,\qquad b+1<(1+\delta)k,\qquad\mu>1.
\end{equation}
These conditions are possible because $r>1/q>1$ and $r<1+\delta$. Fix $\tau\ge0$, let $m_h=\lceil b^{\tau h}\rceil$ and
$G_h=H_h[m_h]$. Then $G_h$ is connected, has
\[
n_h=m_h\Bigl(k+1+\frac{b^{h+1}-1}{b-1}\Bigr)
\]
vertices, so that $\log n_h=(1+\tau)h\log b+O(1)$, and by Lemma~\ref{lem:spectral} and~\eqref{eq:params},
\[
\Delta(G_h)=m_h(b+1)<(1+\delta)m_hk\le(1+\delta)\lambda(G_h),\qquad\lambda(G_h)=\Theta(m_h)=\Theta\bigl(n_h^{\tau/(1+\tau)}\bigr).
\]
Taking $\tau=a/(1-a)$ gives $\lambda(G_h)=\Theta(n_h^a)$. Theorem~\ref{thm:lower} therefore follows from the next
statement.

\begin{proposition}\label{prop:lower}
With $p_h=q/\lambda(G_h)$ and every $0<c<\log\mu/((1+\tau)\log b)$,
\[
\PP\bigl(L_1((G_h)_{p_h})\ge n_h^c\bigr)\to1\qquad(h\to\infty).
\]
\end{proposition}

\subsection{An embedded branching process}
Write $m=m_h$ and $p=p_h$. Fix a vertex $v_0$ in the block of a tree vertex $x_0$, and explore downwards as follows.
Given a selected vertex $v$ in the block of a tree vertex $x$, for each child $y$ of $x$ test the $m$ edges from $v$ to
the block of $y$; if at least one is open, select one of their endpoints by a fixed rule and continue from it. Each
selected vertex lies in the open component of $v_0$, distinct selected vertices lie in distinct blocks, and the tested
edge sets are pairwise disjoint. Given the past, a child block is reached with probability
\begin{equation}\label{eq:success}
1-(1-p)^m\ge1-e^{-mp}=1-e^{-q/\lambda(H_h)}\ge1-e^{-q/(k+1)}=:a_0,
\end{equation}
by Lemma~\ref{lem:spectral}. Thinning each success independently down to probability $a_0$, the selected vertices
contain a Galton--Watson tree, started from $v_0$ and cut at the leaves of the subtree of $x_0$, whose offspring
distribution is $X\sim\mathrm{Bin}(b,a_0)$ with $\EE X=\mu>1$.

\begin{lemma}\label{lem:secondmoment}
Let $Z_\ell$ be the size of generation $\ell$ of a Galton--Watson process with $Z_0=1$ and offspring mean $\mu>1$ and
variance $\sigma^2$. Then $\PP(Z_\ell\ge\mu^\ell/2)\ge s_0$ for all $\ell\ge0$, where
$s_0=\bigl(4+4\sigma^2/(\mu(\mu-1))\bigr)^{-1}$.
\end{lemma}

\begin{proof}
Conditioning on $Z_\ell$ gives $\EE Z_{\ell+1}^2=\mu^2\EE Z_\ell^2+\sigma^2\mu^\ell$, hence
\[
\frac{\EE Z_\ell^2}{\mu^{2\ell}}=1+\frac{\sigma^2}{\mu(\mu-1)}\bigl(1-\mu^{-\ell}\bigr)\le\frac1{4s_0}.
\]
The Paley--Zygmund inequality $\PP(Z\ge\EE Z/2)\ge(\EE Z)^2/(4\EE Z^2)$ (see~\cite[Section~5.5]{LP16}) gives the claim.
\end{proof}

\begin{proof}[Proof of Proposition~\ref{prop:lower}]
Fix $\eta\in(0,1)$ and let $j=\lfloor\eta h\rfloor$. The $b^j$ subtrees of $T_{b,h}$ rooted at level $j$, together
with their blocks, have pairwise disjoint edge sets. Run the exploration above from one vertex in the root block of each
subtree, using no edge to a parent block or to the clique. These explorations are independent, and by
Lemma~\ref{lem:secondmoment} each reaches at least $\mu^{h-j}/2$ distinct vertices at depth $h-j$ of its subtree with
probability at least $s_0$. All vertices reached by one exploration lie in one open component. Therefore
\[
\PP\bigl(L_1((G_h)_p)<\tfrac12\mu^{h-j}\bigr)\le(1-s_0)^{b^j}\to0.
\]
Given $c<\log\mu/((1+\tau)\log b)$, choose $\eta$ with $c<(1-\eta)\log\mu/((1+\tau)\log b)$. Since $\log
n_h=(1+\tau)h\log b+O(1)$, we have $\mu^{h-j}/2\ge n_h^c$ for large $h$.
\end{proof}

\begin{remark}\label{rem:example}
For $\varepsilon=0.02$, $\delta=0.1$, $k=100$ and $b=108$ the conditions~\eqref{eq:params} hold, $\Delta/\lambda\le
109/100$, $\mu=1.04285\ldots$ and $\log\mu/\log b=0.00896\ldots$. With $m_h=108^h$ both $\lambda$ and $\Delta$ are of
order $\sqrt{n}$, and $L_1\ge n^{0.004}$ with high probability. The large component sits inside one subtree, so it does
not use the bridge. Since $k\ge2\sqrt b$ gives $\mu<qb/(k+1)<q\sqrt b/2$, the exponent in
Proposition~\ref{prop:lower} is less than $1/(2(1+\tau))$, while $B=\Theta(n^{1/(2(1+\tau))})$ for these graphs; this
is consistent with~\eqref{eq:cmw}.
\end{remark}

\section{Logarithmic components below the margin}\label{sec:strict}

\begin{proof}[Proof of Proposition~\ref{prop:strict}]
Put $\kappa=1-q(1+\delta)$, so that $p\Delta\le q\Delta/\lambda\le1-\kappa$. Explore the component of a vertex $x$ in
breadth-first order, adding independent trials whenever a vertex has fewer than $\Delta$ unexplored neighbours. The
total progeny is then dominated by that of a Galton--Watson process with offspring $\mathrm{Bin}(\Delta,p)$. If the
total progeny is at least $s$, the first $N=s-1$ particles have at least $N$ children in total, and these children are
a sum $S$ of $N\Delta$ independent Bernoulli$(p)$ variables with $\EE S\le(1-\kappa)N$. With $e^t=1/(1-\kappa)$,
\[
\PP(S\ge N)\le e^{-tN}\EE e^{tS}\le\exp\{-tN+(1-\kappa)N(e^t-1)\}=\exp\{N[\kappa+\log(1-\kappa)]\}\le
e^{-\kappa^2N/2}.
\]
A union bound over $x$ gives $\PP(L_1(G_p)\ge s)\le n\exp(-\kappa^2(s-1)/2)$, which is $o(1)$ for
$s=1+3\kappa^{-2}\log n$.

For the cycle, $\lambda=\Delta=2$ and $p=q/2$. Split its edges into $\lfloor n/\ell\rfloor$ disjoint paths of
$\ell=\lfloor c\log n\rfloor$ edges with $0<c<1/\log(1/p)$. The events that these paths are open are independent with
probability $p^\ell$ each, and $\lfloor n/\ell\rfloor p^\ell\to\infty$, so with high probability one of them occurs and
$L_1\ge\ell+1$.
\end{proof}

\section{The boundary case}\label{sec:boundary}

\begin{proof}[Proof of Theorem~\ref{thm:boundary}]
Assume $p>0$. Let $u=(I-pA)^{-1}\one=\sum_{j\ge0}p^jA^j\one$, which converges because $p\lambda\le q<1$, and let
$U=\max_xu_x$. By~\eqref{eq:walk}, $p^jW_j(x)\le q^jB$, and $p^jW_j(x)\le(p\Delta)^j\le1$. With
$J=\lceil\log B/\log(1/q)\rceil$ the first $J$ terms of $u_x-1$ are at most one and the rest are at most
$q^{j-J}$, so
\begin{equation}\label{eq:U}
U\le1+J+\frac q{1-q}\le1+J+\frac1\varepsilon.
\end{equation}
Since $B\le\sqrt n$, this gives $U=O_\varepsilon(\log n)$.

Consider the multitype branching process with type set $V(G)$ in which a particle of type $x$ has, independently for
each neighbour $y$ of $x$, one child of type $y$ with probability $p$. Exploring the component of $x$ and using fresh
independent trials for edges to vertices already discovered shows that $|C_x|$ is dominated by the total progeny $T$ of
this process started from type $x$. The expected total progeny vector is $u$, and $pAu=u-\one$.

Let $a=1/(2U^2)$ and $\theta=1/(4U^2)$, and let $T_g$ be the progeny in the first $g$ generations. We show by induction
on $g$ that $\EE e^{\theta T_g}\le e^{au_x}$ for a process started at $x$. For $g=0$, $T_0=1$ and $\theta\le a\le
au_x$. For the induction step, independence of the children and $e^z-1\le z+z^2$ for $0\le z\le1$, applied with
$z=au_y\le1/2$, give
\[
\EE e^{\theta T_{g+1}}\le\exp\Bigl[\theta+p\sum_{y\sim x}\bigl(e^{au_y}-1\bigr)\Bigr]
\le\exp\bigl[\theta+a(u_x-1)+a^2U(u_x-1)\bigr]\le e^{au_x},
\]
using $p\sum_{y\sim x}u_y=u_x-1$, $u_y^2\le Uu_y$, $u_x-1\le U$ and $\theta+a^2U^2=a$. By monotone convergence
$\EE e^{\theta T}\le e^{aU}\le e^{1/2}$, so
\[
\PP(|C_x|\ge s)\le\PP(T\ge s)\le\exp\Bigl(\frac12-\frac s{4U^2}\Bigr).
\]
A union bound over $x$ with $s=4U^2(3\log n+1)$ gives $\PP(L_1(G_p)\ge s)\le n^{-2}$, and $s=O_\varepsilon((\log
n)^3)$. If $\Delta\le(1+\delta)\lambda$, $q(1+\delta)=1$ and $p\le q/\lambda$, then $p\Delta\le1$.
\end{proof}

\section{A refinement of the resolvent bound}\label{sec:refine}

\begin{proof}[Proof of Theorem~\ref{thm:refine}]
Let $u=\sum_{j\ge0}(a\lambda)^{-j}A^j\one$. By~\eqref{eq:walk} and $W_j\le\Delta^j$, $(a\lambda)^{-j}W_j(x)\le
a^{-j}\min\{R^j,B\}$, so $1\le u_x\le U_a$. The weights $w=u/U_a$ satisfy $U_a^{-1}\le w_x\le1$ and
$Aw=a\lambda(w-U_a^{-1}\one)\le a\lambda w$.

Explore the component $C$ of a vertex $x$, querying each edge from a discovered vertex to an undiscovered one at most
once. Before the $i$-th query, let $z_i\in[0,1]$ be the weight of the undiscovered endpoint (after the exploration
stops, pad with queries of weight zero), and let $Y_i$ be the indicator that the edge is open. Given the past, $Y_i$ is
Bernoulli$(p)$, so for $t\in[0,1/2]$
\[
\EE\bigl[e^{tY_iz_i}\,\big|\,\mathcal F_{i-1}\bigr]\le\exp\bigl(p(e^{tz_i}-1)\bigr)\le\exp\bigl(pt(1+t)z_i\bigr).
\]
Hence $M_i=\exp\bigl(\sum_{l\le i}(tY_lz_l-pt(1+t)z_l)\bigr)$ is a supermartingale with $M_0=1$. When the exploration
stops, $\sum Y_lz_l=w(C)-w_x\ge w(C)-1$ and $\sum z_l\le\sum_{v\in C}(Aw)_v\le a\lambda w(C)$. With $t=\theta/2$ and
$pa\lambda\le qa=1-\theta$,
\[
\log M\ge\frac\theta2\bigl(w(C)-1\bigr)-\frac\theta2\Bigl(1+\frac\theta2\Bigr)(1-\theta)w(C)\ge\frac{\theta^2}4w(C)-1.
\]
By Markov's inequality $\PP(w(C)\ge z)\le\exp(1-\theta^2z/4)$. Since $|C|\le U_aw(C)$, a union bound over $x$ at
$|C|\ge\lceil4\theta^{-2}U_a(1+3\log n)\rceil$ gives the bound $n\cdot n^{-3}$.
\end{proof}

\begin{remark}\label{rem:exponent}
If $R=1$ then $U_a=a/(a-1)$. If $R>1$, let $J=\lfloor\log B/\log R\rfloor$; then $R^j\le B$ exactly for $j\le J$, and
\[
U_a=1+\sum_{j=1}^{J}\Bigl(\frac Ra\Bigr)^j+\frac{Ba^{-J}}{a-1}.
\]
For fixed $R$ and $a$ this is $O(1)$ if $R<a$, $O(1+\log B)$ if $R=a$, and $O(B^{1-\log a/\log R})$ if $R>a$. If
$R<1/q$, choosing $a\in(R,1/q)$ recovers $L_1=O(\log n)$, in line with Proposition~\ref{prop:strict}. If $R>1/q$ is
fixed and $0<\zeta<-\log q/\log R$, choose $a<1/q$ with $1-\log a/\log R\le\log(qR)/\log R+\zeta$; this gives
\eqref{eq:refine}, whose exponent $\log(qR)/\log R+\zeta$ is less than one. The constants depend on $R$, $a$ and
$\varepsilon$, so~\eqref{eq:refine} is a statement for fixed $R$, not a bound uniform in $R$.
\end{remark}

\section{Examples and computations}\label{sec:examples}

\begin{remark}\label{rem:examples}
The star $K_{1,t}$ has $\lambda=\sqrt t$, $\Delta=t$ and $n=t+1$; at $p=q/\sqrt t$ the centre keeps
$\mathrm{Bin}(t,q/\sqrt t)$ leaves, so $L_1=(q+o(1))\sqrt t$ with high probability, while~\eqref{eq:cmw} gives
$O(\sqrt t\log t)$. Joining $K_{k+1}$ by one edge to a leaf of $K_{1,D}$ with $k+1\le D\le k^2$ gives $\Delta=D$ and
$k\le\lambda\le k+1$, and the largest component is of order $D/k+\log k$; such graphs realise every ratio $\Delta/\lambda$
between about $1$ and $\sqrt\Delta$, but their components grow only through the star. Theorem~\ref{thm:lower} uses
a branching tree instead, which gives polynomial growth at a bounded ratio.
\end{remark}

\begin{remark}\label{rem:sims}
The repository contains a simulation of $1{,}200$ percolation samples on $30$ host graphs of up to $1{,}118{,}688$
vertices, all at $p=(1-\varepsilon)/\lambda$, with $40$ samples per host. The hosts are the balanced complete bipartite
graph, $K_{s,s^3}$, the star with and without a clique, the random graph $G(n,6/(n-1))$ with four added vertices of
degree about $n^{0.6}2^{-i}$, $i=0,1,2,3$ (all with $\varepsilon=0.2$), and the graphs of Section~\ref{sec:lower}
with $k=12$ and $b=16$, with and without blocks of size $2^h$ (with $\varepsilon=0.1$). Least-squares slopes of the
logarithm of the median of $L_1$ against $\log n$ over the tested sizes are $0.17$ and $0.20$ for the two near-regular
families, $0.18$ for the random graph with added vertices, $0.48$ for $K_{s,s^3}$, $0.49$ for the star, $0.52$ for the star with a clique and $0.27$ for
the balanced complete bipartite graph, whose components are logarithmic by Proposition~\ref{prop:strict}. These slopes
are descriptive: the tested heights are small, and the data do not distinguish a small power from slow polylogarithmic
growth. The computations use NumPy~\cite{NumPy} and
SciPy~\cite{SciPy} and are not part of any proof.
\end{remark}

\section*{Acknowledgements and funding}
This work was done by multiple AI agents, Claude (Anthropic) and Codex (OpenAI), connected to Math
Brain~\cite{MathBrainRepo} and working together inside the Brain framework~\cite{BrainFramework}, a shared, persistent
memory in which the agents record claims, proofs, failed attempts and methods, retrieve one another's results,
coordinate who works on what, and audit each other's work. For more on the details behind the Brain and Math Brain, see
the Math Brain paper~\cite{MathBrainPaper}, the Math Brain repository~\cite{MathBrainRepo} and the Brain framework
repository~\cite{BrainFramework}.

The research presented in this paper was supported by the European Research Council (ERC) under the European Union's
Horizon Europe research and innovation programme (grant agreement No.~101041711), by the Simons Foundation as part of
the Collaboration on the Mathematical and Scientific Foundations of Deep Learning, by Heights Labs, by Convex Nexus
Capital, by the Israel Science Foundation (grant number 2258/19), and by the Israel Science Foundation (ISF Grant
4101/25). Computational resources were provided in part by Google Cloud through its education credits programme (credit
530550103).

\section*{Declaration of competing interest}
The author holds equity in Heights Labs and in Convex Nexus Capital and has a management role at Convex Nexus Capital.
These entities are included in the funding declaration above.

\section*{Declaration of generative AI and AI-assisted technologies}
A Codex agent (OpenAI) proved the results and wrote the code of Remark~\ref{rem:sims}; a second Codex agent checked the
proofs and the literature as a referee; Claude agents (Anthropic) checked the steps and prepared the manuscript. The
author directed the work and is responsible for the content of the paper and for decisions concerning its
dissemination.

\section*{Code availability}
The code and data of Remark~\ref{rem:sims} are in the repository
\begin{center}\url{https://github.com/yspennstate/near-regular-percolation-components}.\end{center}

\appendix
\section{Notation}\label{app:notation}

\begin{center}
\begin{tabular}{p{3.0cm}p{8.0cm}p{3.1cm}}
\toprule
Symbol & Meaning & First appears\\
\midrule
$G$, $n$, $A$ & finite simple graph, its order, its adjacency matrix & Section~\ref{sec:intro}\\
$\lambda$, $\Delta$, $d(x)$ & spectral radius, maximum degree, degree of $x$ & Section~\ref{sec:intro}\\
$G_p$, $L_1$ & percolated graph, order of its largest component & Section~\ref{sec:intro}\\
$\varepsilon$, $q$ & subcritical margin, $q=1-\varepsilon$ & Section~\ref{sec:intro}\\
$\delta$ & near-regularity parameter, $\Delta\le(1+\delta)\lambda$ & Section~\ref{sec:intro}\\
$R$, $B$ & $\Delta/\lambda$ and $\sqrt{n\Delta}/\lambda$ & Section~\ref{sec:intro}\\
$c$ & exponent of the large component & Theorem~\ref{thm:lower}\\
$\one$, $e_x$ & all-ones vector, basis vector at $x$ & Section~\ref{sec:intro}\\
$\PP$, $\EE$ & probability, expectation & Section~\ref{sec:intro}\\
$v_0$, $x_0$ & starting vertex, its tree vertex & Section~\ref{sec:lower}\\
$a$, $\theta$, $U_a$ & resolvent parameter, $1-qa$, weight bound & Theorem~\ref{thm:refine}\\
$\zeta$ & slack in the exponent & \eqref{eq:refine}\\
$W_k(x)$ & number of walks of length $k$ from $x$ & \eqref{eq:walk}\\
$T_{b,h}$, $H_h$ & $b$-ary tree of height $h$, tree joined to $K_{k+1}$ & Section~\ref{sec:lower}\\
$H_h[m]$, $J_m$ & blow-up with blocks of size $m$, all-ones matrix & Section~\ref{sec:lower}\\
$r$, $\mu$, $\tau$ & ratio $b/k$, offspring mean, blow-up exponent & Section~\ref{sec:lower}\\
$m_h$, $G_h$, $n_h$, $p_h$ & block size, constructed graph, its order, $q/\lambda(G_h)$ & Section~\ref{sec:lower}\\
$a_0$, $X$ & success bound, offspring variable & \eqref{eq:success}\\
$Z_\ell$, $\sigma^2$, $s_0$ & generation size, offspring variance, success bound & Lemma~\ref{lem:secondmoment}\\
$\eta$, $j$ & fraction of the height, level of the subtrees & Proposition~\ref{prop:lower}\\
$\kappa$, $S$, $N$ & $1-q(1+\delta)$, sum of offspring, its length & Section~\ref{sec:strict}\\
$u$, $U$, $J$ & resolvent vector, its maximum, cut-off index & Section~\ref{sec:boundary}\\
$T$, $T_g$, $C_x$ & total progeny, through generation $g$, component of $x$ & Section~\ref{sec:boundary}\\
$w$, $z_i$, $Y_i$, $M_i$ & weights, query weight, query indicator, supermartingale & Section~\ref{sec:refine}\\
\bottomrule
\end{tabular}
\end{center}

\begin{thebibliography}{99}
\bibitem{ABS04} N.~Alon, I.~Benjamini and A.~Stacey, Percolation on finite graphs and isoperimetric inequalities,
\emph{Ann.\ Probab.} \textbf{32} (2004), no.~3, 1727--1745. \doi{10.1214/009117904000000414}
\bibitem{BBCR10} B.~Bollobás, C.~Borgs, J.~Chayes and O.~Riordan, Percolation on dense graph sequences, \emph{Ann.\
Probab.} \textbf{38} (2010), no.~1, 150--183. \doi{10.1214/09-AOP478}
\bibitem{CMW26} M.~Christoph, P.~Morawski and Y.~Wigderson, The critical probability for percolation on finite graphs,
preprint (2026), \arxiv{2608.19145}.
\bibitem{DEKK25} S.~Diskin, J.~Erde, M.~Kang and M.~Krivelevich, Percolation on high-dimensional product graphs,
\emph{Random Structures Algorithms} \textbf{66} (2025), no.~1, e21268. \doi{10.1002/rsa.21268}; \arxiv{2209.03722}.
\bibitem{NumPy} C.~R.~Harris et al., Array programming with NumPy, \emph{Nature} \textbf{585} (2020), 357--362.
\doi{10.1038/s41586-020-2649-2}
\bibitem{LP16} R.~Lyons and Y.~Peres, \emph{Probability on Trees and Networks}, Cambridge University Press, New York,
2016. \doi{10.1017/9781316672815}
\bibitem{MathBrainPaper} Y.~Shmalo, \emph{Math Brain: shared, auditable memory for multi-agent mathematical research},
October 2026, \url{https://github.com/yspennstate/math-brain-public/tree/main/papers/math-brain-2026-10-24}.
\bibitem{MathBrainRepo} Y.~Shmalo, Math Brain: public snapshot, code and tests, 2026,
\url{https://github.com/yspennstate/math-brain-public}.
\bibitem{BrainFramework} Y.~Shmalo, Brain Framework, public software repository, 2026,
\url{https://github.com/yspennstate/brain-framework-public}.
\bibitem{SciPy} P.~Virtanen et al., SciPy 1.0: fundamental algorithms for scientific computing in Python, \emph{Nature
Methods} \textbf{17} (2020), 261--272. \doi{10.1038/s41592-019-0686-2}
\end{thebibliography}
\end{document}
